\section{Diffraction Model}
\label{sec:diffraction_model}
We refer to the model as a forward model because it follows the physical sequence from ray to sample to detector. The incident beam on the sample is represented by weighted source states sampled in position, direction, and wavelength. Each incident ray is described by its incident angles $(\theta_i,\phi_i)$ and wavevector magnitude $|\mathbf{k}_i|=2\pi/\lambda$, so its state is written as $\mathbf{k}_i(|\mathbf{k}_i|,\theta_i,\phi_i)$. Sample composition and crystallinity determine a reciprocal-space structure function $S(\mathbf{k}_i,\mathbf Q)$. Restricting this density to the Ewald sphere defined by $\mathbf{k}_i$ produces an Ewald surface decorated by the intensities in $\mathbf{Q}$. The corresponding $\mathbf{k}_f$ vectors carry the selected intensities toward the shared geometry of the chosen X-ray collection method, namely a 2D area detector.

\subsection{Scattering and detector geometry}
\label{sec:scattering_detector_geometry}
For one incident state, let $\mathbf{k}_i$ denote the incident wavevector and $\mathbf{k}_f$ the scattered wavevector. For wavelength $\lambda$, elastic scattering preserves their magnitudes,
\begin{equation}
  |\mathbf{k}_i|=|\mathbf{k}_f|=\frac{2\pi}{\lambda}.
  \label{eq:elastic_scattering}
\end{equation}
The momentum transfer is
\begin{equation}
  \mathbf Q=\mathbf{k}_f-\mathbf{k}_i.
  \label{eq:momentum_transfer}
\end{equation}
For an infinite, perfectly ordered crystal, an allowed reflection $(h,k,\ell)$ supplies
\begin{equation}
  \mathbf G_{hk\ell}=h\mathbf a^{*}+k\mathbf b^{*}+\ell\mathbf c^{*},
  \qquad \mathbf Q=\mathbf G_{hk\ell},
  \label{eq:bragg_matching}
\end{equation}
where $\mathbf a^{*}$, $\mathbf b^{*}$, and $\mathbf c^{*}$ include the factor $2\pi$. The corresponding exit wavevector is
\begin{equation}
  \mathbf{k}_f=\mathbf{k}_i+\mathbf G_{hk\ell}.
  \label{eq:exit_wavevector}
\end{equation}
Equations~\eqref{eq:bragg_matching} and~\eqref{eq:exit_wavevector} define the ideal integer Bragg-point limit. Replacing the integer order $\ell$ by a continuous coordinate $L$ generalizes the same construction to a reciprocal rod, with the Bragg orders lying at $L=\ell$.

Figure~\ref{fig:waxs_geometry_system} shows how an accepted exit direction reaches the area detector. In the laboratory frame, $y$ follows the nominal incident beam, $z$ is vertical, and $x$ is transverse. The scattering angle $2\theta$ is the angle between $\mathbf{k}_i$ and $\mathbf{k}_f$. The outgoing-ray azimuth $\phi_f$ is measured about the incident direction, with zero toward $+z$ in the reference geometry.

The detector axes $x'$ and $z'$ lie in its plane, and $\hat{\mathbf n}_{\mathrm{det}}$ points out of its active face toward the sample. The nominal sample-to-detector distance $D$ is measured along the beam axis in the untilted reference geometry. Detector tilts $\beta_D$ and $\gamma_D$ change the plane orientation. Supporting Information Sec.~S1 gives the corresponding detector-frame coordinates, signed angular convention, and ray--plane intersection equations.

\begin{figure}[H]
  \centering
  \resizebox{0.6\linewidth}{!}{% Detector geometry TikZ figure (to be \input into main document)

\tdplotsetmaincoords{70}{120}
\begin{tikzpicture}[tdplot_main_coords, scale=1.5]

    %=========================
    % Parameters
    %=========================
    % Geometry angles (keep detector unrotated)
    \pgfmathsetmacro{\alpha}{0}         % detector tilt about x (deg)
    \pgfmathsetmacro{\betaang}{0}       % detector rotation about z (deg)

    % Display-only angles (nonzero annotations)
    \pgfmathsetmacro{\betaVis}{-20}      % shown beta angle (deg)
    \pgfmathsetmacro{\gammaVis}{20}     % shown gamma angle (deg)

    \pgfmathsetmacro{\detW}{2}           % detector width
    \pgfmathsetmacro{\detH}{2}           % detector height
    \pgfmathsetmacro{\detT}{0.2}         % detector thickness

    \pgfmathsetmacro{\detHw}{0.5*\detW}  % half-width
    \pgfmathsetmacro{\detHh}{0.5*\detH}  % half-height

    % Geometry sines and cosines
    \pgfmathsetmacro{\sina}{sin(\alpha)}
    \pgfmathsetmacro{\cosa}{cos(\alpha)}
    \pgfmathsetmacro{\sinb}{sin(\betaang)}
    \pgfmathsetmacro{\cosb}{cos(\betaang)}

    % Display trig
    \pgfmathsetmacro{\sinBetaVis}{sin(\betaVis)}
    \pgfmathsetmacro{\cosBetaVis}{cos(\betaVis)}
    \pgfmathsetmacro{\sinGammaVis}{sin(\gammaVis)}
    \pgfmathsetmacro{\cosGammaVis}{cos(\gammaVis)}

    % Detector center (front face)
    \coordinate (detCenter) at (-3,0,0);

    %=========================
    % Detector local basis (tilt about x then rotation about z)
    %=========================
    % Horizontal (width) direction
    \coordinate (detU) at ({-\detHw*\sinb},{\detHw*\cosb},0);

    % Vertical (height) direction
    \coordinate (detV) at ({\detHh*\sina*\cosb},
                           {\detHh*\sina*\sinb},
                           {\detHh*\cosa});

    % Thickness direction (front to back)
    \coordinate (detTvec) at ({-\detT*\cosa*\cosb},
                              {-\detT*\cosa*\sinb},
                              {\detT*\sina});

    % Components of detV for use in arc centers
    \pgfmathsetmacro{\Vx}{\detHh*\sina*\cosb}
    \pgfmathsetmacro{\Vy}{\detHh*\sina*\sinb}
    \pgfmathsetmacro{\Vz}{\detHh*\cosa}

    % Bottom center coordinates: detFBC = detCenter - detV
    \pgfmathsetmacro{\FBCx}{-3 - \Vx}
    \pgfmathsetmacro{\FBCy}{0  - \Vy}
    \pgfmathsetmacro{\FBCz}{0  - \Vz}

    % Bottom right coordinates: detFBR = detCenter + detU - detV
    \pgfmathsetmacro{\FBRx}{-3 - \detHw*\sinb - \Vx}
    \pgfmathsetmacro{\FBRy}{    \detHw*\cosb - \Vy}
    \pgfmathsetmacro{\FBRz}{          0      - \Vz}

    % Top right coordinates: detFTR = detFBR + 2*detV
    \pgfmathsetmacro{\FTRx}{\FBRx + 2*\Vx}
    \pgfmathsetmacro{\FTRy}{\FBRy + 2*\Vy}
    \pgfmathsetmacro{\FTRz}{\FBRz + 2*\Vz}

    %=========================
    % Detector bottom-center point
    %=========================
    \coordinate (detFBC) at ({\FBCx},{\FBCy},{\FBCz});

    %=========================
    % Detector geometry
    %=========================
    % Front face corners
    \coordinate (detFBL) at ($ (detCenter) - (detU) - (detV) $); % bottom left
    \coordinate (detFBR) at ($ (detCenter) + (detU) - (detV) $); % bottom right
    \coordinate (detFTL) at ($ (detCenter) - (detU) + (detV) $); % top left
    \coordinate (detFTR) at ($ (detCenter) + (detU) + (detV) $); % top right

    % Back face corners
    \coordinate (detBBL) at ($ (detFBL) + (detTvec) $);
    \coordinate (detBBR) at ($ (detFBR) + (detTvec) $);
    \coordinate (detBTL) at ($ (detFTL) + (detTvec) $);
    \coordinate (detBTR) at ($ (detFTR) + (detTvec) $);

    % Hidden-edge style
    \tikzset{
      detectorHidden/.style={
        line width=0.45pt,
        dashed,
        dash pattern=on 2pt off 1.2pt
      }
    }

    % Draw back face edge-by-edge
    \draw[detectorHidden] (detBBL) -- (detBBR); % bottom back
    \draw[very thick]     (detBBR) -- (detBTR); % right back
    \draw[very thick]     (detBTR) -- (detBTL); % top back
    \draw[detectorHidden] (detBTL) -- (detBBL); % left back

    % Draw front face
    \draw[very thick]
      (detFBL) -- (detFBR) -- (detFTR) -- (detFTL) -- cycle;

    % Connector edges (thickness)
    \draw[very thick]     (detFBR) -- (detBBR);
    \draw[very thick]     (detFTR) -- (detBTR);
    \draw[very thick]     (detFTL) -- (detBTL);
    \draw[detectorHidden] (detFBL) -- (detBBL);

    % Detector label
    \node at ($(detCenter)+(0,0,1.3)$) {\scriptsize Detector};
    %=========================
    % Detector-frame axes (schematic)
    %=========================
    \draw[->, black, line width=0.2mm] (detCenter) -- ($ (detCenter) + 0.9*(detU) $)
      node[anchor=west, scale=0.5, xshift = -12, yshift = -6] {$x'$};
    \draw[->, black, line width=0.2mm] (detCenter) -- ($ (detCenter) + 0.9*(detV) $)
      node[anchor=south, scale=0.5, xshift = -8, yshift = -20] {$z'$};


% Detector normal anchor point on the front face
\coordinate (detNormalBase) at ($ (detCenter) + 0.55*(detU) - 0.35*(detV) $);
\coordinate (detNormalTip)  at ($ (detNormalBase) - 2.6*(detTvec) $);

% Optional style definitions
\tikzset{
  detNormalArrow/.style={
    ->,
    black,
    line width=0.3mm,
    preaction={draw, white, line width=0.9mm}
  },
  detNormalPoint/.style={
    fill=black,
    draw=white,
    line width=0.25pt
  }
}

% Mark the point on the detector surface
\filldraw[detNormalPoint] (detNormalBase) circle (0.7pt);

% Draw the detector normal with a subtle white halo
\draw[detNormalArrow] (detNormalBase) -- (detNormalTip)
  node[pos=0.92, anchor=north west, scale=0.52,
       xshift=2pt, yshift=-1pt,
       fill=white, inner sep=1pt, text=black]
  {$\hat{n}_{\mathrm{det}}$};
  
    %=========================
    % Beta arc, line, and label at top-right (display-only)
    %=========================
    \pgfmathsetmacro{\rBetaArc}{0.3}
    \pgfmathsetmacro{\betastart}{min(0,\betaVis)}
    \pgfmathsetmacro{\betaend}{max(0,\betaVis)}
    \pgfmathsetmacro{\betamid}{0.5*(\betastart+\betaend)}

    % Dashed line for beta from top-right in a local xz plane
    \coordinate (betaBase) at ({\FTRx},{\FTRy},{\FTRz});
    \coordinate (betaDet) at ({\FTRx + \rBetaArc*\sinBetaVis},
                              {\FTRy},
                              {\FTRz + \rBetaArc*\cosBetaVis});
    \draw[orange, dash pattern=on 1pt off 0.7pt] (betaBase) -- (betaDet);

    % Vertical z bar at top-right
    \coordinate (detFTRzBar) at ($ (detFTR) + (0,0,0.3) $);
    \draw[orange,thick] (detFTR) -- (detFTRzBar);

    % Arc for beta at top-right
    \draw[orange, dash pattern=on 1pt off 0.7pt, line width=0.1mm] 
      plot[domain=\betastart:\betaend, samples=30, variable=\t]
        ({\FTRx + \rBetaArc*sin(\t)},
         {\FTRy},
         {\FTRz + \rBetaArc*cos(\t)});

    % Label for beta at top-right
    \coordinate (betaLabel) at ({\FTRx + 1.1*\rBetaArc*sin(\betamid)},
                                {\FTRy},
                                {\FTRz + 1.1*\rBetaArc*cos(\betamid)});
    \node[orange, scale=0.6, anchor=south east,
          xshift=6pt, yshift=-5pt] at (betaLabel) {$\beta$};

    %=========================
    % Gamma arc, lines, and label at top-right (display-only)
    %=========================
    \pgfmathsetmacro{\rArc}{0.4}
    \pgfmathsetmacro{\alphab}{90+\gammaVis}   % rotated edge direction (display)
    \pgfmathsetmacro{\alphax}{90}             % reference +y
    \pgfmathsetmacro{\alphastart}{min(\alphab,\alphax)}
    \pgfmathsetmacro{\alphaend}{max(\alphab,\alphax)}
    \pgfmathsetmacro{\alphamid}{0.5*(\alphastart+\alphaend)}

    \coordinate (gammaBase) at ({\FTRx},{\FTRy},{\FTRz});
    % display edge direction using (-sin, cos)
    \coordinate (gammaWidthEnd) at ({\FTRx - 0.4*\sinGammaVis},
                                    {\FTRy + 0.4*\cosGammaVis},
                                    {\FTRz});
    % reference +y direction
    \coordinate (gammaXEnd) at ({\FTRx},
                                {\FTRy + 0.4},
                                {\FTRz});
    \draw[oiPurple,thick] (gammaBase) -- (gammaWidthEnd);
    \draw[oiPurple,thin]  (gammaBase) -- (gammaXEnd);

    % Gamma arc centered at top-right
    \draw[oiPurple, line width=0.1mm, dash pattern=on 1pt off 0.7pt]
      plot[domain=\alphastart:\alphaend, samples=30, variable=\t]
        ({\FTRx + \rArc*cos(\t)},
         {\FTRy + \rArc*sin(\t)},
         {\FTRz});

    % Label for gamma
    \coordinate (gammaLabel) at ({\FTRx + 0.7*cos(\alphamid)},
                                 {\FTRy + 0.7*sin(\alphamid)},
                                 {\FTRz});
    \node[oiPurple, scale=0.6, anchor=south west,
          xshift=-18pt, yshift=-5pt] at (gammaLabel) {$\gamma$};

    %=========================
    % Sample and beams
    %=========================
    % Incident beam from sample to detector center
    \draw[black, dash dot] (-0.5,0,0) -- (detCenter);

    % Sample plane
    \draw[fill=gray]
      (-0.312, -0.2, -0.068) --
      (-0.312,  0.2, -0.068) --
      (-0.688,  0.2,  0.068) --
      (-0.688, -0.2,  0.068) -- cycle;

    \node[scale=0.7, xshift = -15, yshift = 25] at (-0.212, 0.3, -0.088) {\tiny Sample};

    % Incident beam from source to sample
    \draw[black, dash dot] (0.5,0,0) -- (-0.5,0,0);

    %=========================
    % Scattered point and components
    %=========================
    \pgfmathsetmacro{\offY}{0.8}  % along U0 (width)
    \pgfmathsetmacro{\offZ}{0.3}  % along V0 (height)

    \coordinate (scatP) at ({-3 + \offZ*\sina*\cosb - \offY*\sinb},
                            {\offZ*\sina*\sinb + \offY*\cosb},
                            {\offZ*\cosa});

    % Width-component point: detCenter + offY * u_hat, u_hat = (-sinb, cosb, 0)
    \coordinate (scatPwidth) at ({-3 - \offY*\sinb},
                                 {0  + \offY*\cosb},
                                 {0});

    \draw[dash pattern=on 2pt off 1pt, green] (detCenter) -- (scatP);
    \draw[dash pattern=on 2pt off 1pt,green] (detCenter) -- (scatPwidth);
    \draw[dash pattern=on 2pt off 1pt,green] (scatPwidth) -- (scatP);

    % Scattered beam vector r
    \draw[red,->] (-0.5,0,0) -- (scatP);
    \node[scale=0.7, red, xshift=12pt, yshift=-2pt]
      at ($ (-0.5,0,0)!0.6!(scatP) $) {\tiny $\vec{k}_f$};

    % Incident vector k_i
    \draw[blue,->] (0.5,0,0) -- (-0.5,0,0);
    \node[scale=0.7, blue, yshift= -0.5mm] at (0.0,0.12,-0.01) {\tiny $\vec{k}_i$};

    %=========================
    % Angle arcs for phi and 2theta
    %=========================
    \tdplotsetrotatedcoords{0}{90}{90}
    \tdplotdrawarc[dashed,tdplot_rotated_coords, cyan]
      {(0,0,-3)}{0.8}{20}{90}{anchor=north west, xshift = -8pt}{\tiny $\phi$}

    \tdplotsetrotatedcoords{90}{-20.56}{0}
    \tdplotdrawarc[dashed, green,tdplot_rotated_coords, line width=0.1mm]
      {(0,1.4,0)}{0.5}{15}{90}
      {anchor=south east}{\tiny $2\theta$}

    %=========================
    % Axes legend
    %=========================
    \draw[->] (.5,0,0) -- (0.3,0,0)
      node[anchor=south west, xshift=-1.5mm, yshift=-0.8mm] {\tiny $y$};
    \draw[->] (.5,0,0) -- (.5,0.2,0) node[anchor=west] {\scriptsize $x$};
    \draw[->] (.5,0,0) -- (.5,0,0.2) node[anchor=south] {\scriptsize $z$};

\end{tikzpicture}
}
  \caption{Laboratory and detector geometry in the untilted reference pose. The blue incident wavevector $\mathbf{k}_i$ reaches the sample, and the red scattered wavevector $\mathbf{k}_f$ reaches the detector. The dash-dotted line continues the incident direction and defines the reference for $2\theta$. The azimuth $\phi_f$ is measured from $+z'$ about that direction and is negative for the ray shown. Dashed detector-plane lines resolve the hit position along $x'$ and $z'$. The laboratory axes are $x,y,z$, the outward detector normal is $\hat{\mathbf n}_{\mathrm{det}}$, and $D$ marks the nominal sample-to-detector distance.}
  \label{fig:waxs_geometry_system}
\end{figure}

\begingroup
\clubpenalty=10000
Source position and propagation angles arising from divergence are sampled from the measured Gaussian beam distributions, while wavelength is sampled from the input wavelength spectrum. Supporting Information Sec.~S1 defines the weighted source-state construction. The external ray intersects the measured sample geometry, and refraction gives the internal vector $\mathbf{k}'_i$. Accepted exit rays must intersect the front face within the active detector bounds. Sample acceptance and detector acceptance are therefore consecutive geometric restrictions.
\par\endgroup

\subsection{Structure Factor}
\label{sec:structure_factor}
For the hexagonal basal plane, every integer in-plane pair $(h,k)$ belongs to the radial family
\begin{equation}
  r\equiv h^2+hk+k^2,
  \qquad
  Q_R(r)=\frac{4\pi}{\sqrt{3}\,a}\sqrt r.
  \label{eq:hexagonal_family_index}
\end{equation}
Members with the same $r$ have the same in-plane radius. Their amplitudes are evaluated individually, and their intensity contributions are summed to represent scattering from the radial family. The family label is distinct from crystallographic multiplicity, and the explicit sum accounts for multiplicity.

For one $(h,k)$ member, the crystallite-frame reciprocal rod is
\begin{equation}
  \mathbf G_{hk}(L)=h\mathbf a^{*}+k\mathbf b^{*}+L\mathbf c^{*}.
  \label{eq:continuous_reciprocal_rod}
\end{equation}
The indices $h$ and $k$ are integers, while $L$ is continuous. An integer Bragg order $(h,k,\ell)$ lies at $L=\ell$. The wavelength-dependent amplitude along the rod is
\begin{equation}
  F_{hk}(L,\lambda_s)=
  \sum_a o_a f_a\!\left(|\mathbf G_{hk}(L)|,\lambda_s\right)
  T_a\!\left(\mathbf G_{hk}(L)\right)
  \exp\!\left[i\mathbf G_{hk}(L)\cdot\mathbf r_a\right].
  \label{eq:continuous_structure_factor}
\end{equation}
Here $o_a$, $f_a$, $T_a$, and $\mathbf r_a$ are the occupancy, atomic scattering factor, atomic displacement (Debye--Waller) factor, and Cartesian site position.

The complex amplitude $F_{hk}$ is the structural input evaluated along the reference rod. The remaining calculation uses the continuous rod intensity $S_{hk}(L,\lambda_s)$. For an ordered finite film, $S_{hk}$ is constructed from $|F_{hk}|^2$. $S_{hk}(L,\lambda_s)\,dL$ is the integrated structural intensity carried by a small segment of the unrotated $(h,k)$ rod. Supporting Information Sec.~S2 gives the finite-stack construction, displacement conventions, and conversion between densities per $dL$ and per dimensional rod length.

The cylinder construction follows directly from the definition of a 2D powder in Fig.~\ref{fig:intro_2x3}(f--g): the crystallites share a mean film normal but are random in in-plane azimuth. If $\mathbf C$ maps the mean crystallite frame into the sample frame, rotating a point on the rod through azimuth $\beta$ gives
\begin{equation}
  \mathbf Q^{\mathrm{cyl}}_{hk}(\beta,L)
  =\mathbf C\,\mathbf R_{\widehat{\mathbf c}^{\,*}}(\beta)\mathbf G_{hk}(L),
  \qquad 0\leq\beta<2\pi.
  \label{eq:powder_cylinder}
\end{equation}
Thus, every off-axis rod sweeps out a cylinder about the film-normal reciprocal-space axis, while the axial rod remains on the normal axis. Members with the same $r$ share the same geometric cylinder, although their amplitudes are evaluated separately. For off-axis cylinders with $Q_R>0$, distributing a density per dimensional rod length around the circumference introduces $1/[2\pi Q_R(r)]$. For the density per $dL$ used here, the physical cylinder-area conversion additionally contains $1/|\mathbf c^*|$ (Supporting Information Sec.~S2). This construction produces the paired off-specular trajectories in Fig.~\ref{fig:intro_2x3}(h) while leaving axial reciprocal points on the common normal.
